\documentclass[aspectratio=169]{beamer}

\usepackage[utf8]{inputenc}
\usepackage{graphicx}
\usepackage{xcolor}
\usepackage[none]{hyphenat}

\usetheme{default}

% Remove navigation symbols
\setbeamertemplate{navigation symbols}{}

\title{Mini Assignment: Deconstructing HCI \& Computer Graphics}

\date{}

\begin{document}

%================================================
% SLIDE 1 - HCI FOCUS
%================================================

\begin{frame}{HCI Focus -- Google Earth}

\begin{center}

\includegraphics[
    width=0.92\paperwidth,
    height=0.60\paperheight,
    keepaspectratio
]{google_earth_hci.png}

\end{center}

\vspace{-0.1cm}

\small

\textbf{Reflection 1: How does the Computer Graphics quality help or hinder
the user's ability to interact with the software?}

High-quality graphics such as realistic 3D views, lighting and shadows help
users understand locations, depth and scale, making Google Earth easier and
more engaging to explore. However, detailed graphics may require more
processing power and can make the software less responsive on lower-end
devices.

\end{frame}


%================================================
% SLIDE 2 - COMPUTER GRAPHICS FOCUS
%================================================

\begin{frame}{Computer Graphics Focus -- Google Earth}

\begin{center}

\includegraphics[
    width=0.92\paperwidth,
    height=0.60\paperheight,
    keepaspectratio
]{google_earth_cg.png}

\end{center}

\vspace{-0.1cm}

\small

\textbf{Reflection 2: What is one simple HCI change you would make to improve
the user interface?}

I would add clearer labels or tooltips for the main navigation controls,
such as the compass, 2D/3D button and zoom controls. This would help new
users understand what each control does and make Google Earth more
beginner-friendly and easier to use.

\end{frame}


\end{document}